\section{Related Work}\label{related_work}

\subsection{Vision-Language Models}

Vision-language models combine visual and textual representations to support multimodal tasks.
These tasks include image captioning and VQA. Models such as BLIP
demonstrated that vision-language pretraining can transfer effectively across both understanding
and generation tasks, including VQA \parencite{li2022blip}. Architecturally, BLIP employs a vision
encoder, a multimodal question encoder and an answer decoder, using cross-attention to integrate visual
and textual representations \parencite{li2022blip}. More recently, generative VLMs instead connect 
visual encoders to autoregressive language models through learned projection mechanisms. SmolVLM combines 
a compact SigLIP-based vision encoder with SmolLM2, emphasising efficient visual-token representations \parencite{marafioti2025smolvlm}.
Similarly, MiniCPM-V 4.6 integrates SigLIP2 with Qwen 3.5 and employs visual-token compression to balance
efficiency and visual understanding \parencite{marafioti2025smolvlm}. These architectural differences
provide a basis for investigating whether post-training effectiveness varies across VLMs.

% VQA is an interesting test of 
% multimodal reasoning since models would need to identify relevant visual evidence and the relate
% it to a natural-language question. The GQA data enhances this setting towards compositional visual
% reasoning. This is done by constructing questions from Visual Genome scene graphs that come from
% functional programs that represent their semantics \parencite{hudson2019gqa}. These structured
% annotations also provide explicit grounding between questions and image objects, making GQA
% suitable for constructing privileged visual regions during training.

\subsection{Knowledge Distillation and On-Policy Learning}

Knowledge distillation is a technique in which information is transferred from a stronger teacher model
to a smaller student by encouraging the student to match the teacher's output distribution
\parencite{hinton2015distilling}. In language generation, however, distillation from a fixed
teacher-generated trajectory can expose the student to sequences produced by the teacher rather
than to states encountered under its own policy. On-policy distillation addresses this
mismatch by allowing the student to generate its own trajectories and obtain dense teacher
supervision along those trajectories.

%% more on results and how cantdirectly translate to vision

\textcite{zhao2601self} extended this idea through On-Policy Self-Distillation (OPSD). Instead
of requiring a separate teacher model, OPSD uses the same model for the teacher and the student under
different contexts. The student receives only the information available during inference, while
the teacher receives additional privileged information. The student generates the trajectory and
training minimises token-level divergence between the student and teacher distributions along
that trajectory. Their experiments focused on mathematical reasoning with large language models.
Their results demonstrated improvements over off-policy distillation methods \parencite{zhao2601self}. However, these findings
were established on mathematical reasoning tasks, where privileged information can provide explicit
guidance towards the correct solution. In multimodal settings, additional visual context does not
necessarily provide an equivalent advantage, leaving the effectiveness of OPSD for VQA less certain.

\subsection{Multimodal On-Policy Self-Distillation}

More recently, \textcite{yuan2026vision} proposed Vision-OPD, extending OPSD to fine-grained 
multimodal understanding. They identified a regional-to-global perception gap in which
multimodal models can answer questions correctly when presented with an evidence-centred image 
crop but fail when viewing the corresponding full image. Vision-OPD therefore uses a crop-conditioned
instance of the model as the teacher and a full-image instance as the student. The student
generates trajectories while token-level supervision is provided by the privileged crop-conditioned
teacher. Their experiments demonstrated substantial improvements in fine-grained visual understanding.
For Qwen3.5-9B, Vision-OPD improves accuracy on V* Bench from 82.72\% to 94.76\% and on ZoomBench from 
52.07\% to 65.80\% \parencite{yuan2026vision}. However, their training examples were specifically
constructed around localised visual evidence, which may limit the transferability of these findings to 
existing VQA datasets.

Although, subsequent work has further highlighted that crop-conditioned supervision is not uniformly 
beneficial. \textcite{hu2026evidencealigned} argue that cropping may actually remove useful
surrounding context. They proposed Evidence-Aligned Distillation (EAD), which selectively weights teach 
corrections according to their alignment with relevant visual evidence. Their method retained approximately 6\% of the
original dense supervision mass while outperforming the compared approaches, suggesting that the relevance of teacher 
supervision is more important than its quantity \parencite{hu2026evidencealigned}. Their findings motivate further evaluation of when privileged visual context
provides useful self-distillation signals.

The present work builds on this emerging research area but focuses specifically on VQA. It evaluates 
OPSD against conventional SFT across three very different VLM architectures. Rather than assuming
that privileged crops necessarily produce a superior teacher, the experiments explicitly measure the 
teacher advantage and examine whether the relative benefit of OPSD varies across model architectures.